\documentclass[10pt,a4paper]{amsart}
\usepackage[margin=19mm]{geometry}
\usepackage{amsmath,amssymb,mathtools}
\usepackage[scr=rsfs,cal=euler]{mathalfa}
\usepackage{xfrac}
\usepackage[unicode,hidelinks,bookmarks=false]{hyperref}
\newcommand{\ie}{i.e.\ }
\newcommand{\cf}{cf.\ }
\newcommand{\resp}{respectively\ }
\newcommand{\Subsection}{\subsection}
\providecommand{\nobreakdash}{\nobreak\hbox{-}}
\setlength{\emergencystretch}{2em}
\multlinegap0pt
\usepackage{amssymb}
\usepackage{mathtools}
\newtheorem{theorem}{Theorem}
\newtheorem{lemma}{Lemma}
\title{Who Pays for a Connected Public Good?}
\author{Marco Tulio Angulo}
\date{Source: arXiv:2609.10565v1 (CC BY 4.0)}
\begin{document}
\maketitle

\noindent\emph{Source and licence.} Who Pays for a Connected Public Good?. Version arXiv:2609.10565v1 (CC BY 4.0), distributed by arXiv under the Creative Commons Attribution 4.0 International licence, http://creativecommons.org/licenses/by/4.0/, as recorded in the arXiv licence metadata for this exact version and shown on the abstract page https://arxiv.org/abs/2609.10565v1. Full licence notice: \texttt{licenses/arxiv-2609.10565-CC-BY-4.0.txt}.\par\smallskip

\noindent\textit{Editorial scope.} Counted unit: the article's theorem thm:graph-bound (source0.tex lines 479--498). The article's complete original proof of it is delivered with it (source0.tex lines 912--1034, one \texttt{proof} environment of 123 lines, which the article places away from the statement). No mathematics is written, repaired or re-derived: the statement and the proof are the article's own lines, in the article's own notation, and no retained line is substituted or re-flowed.  Retained uncounted as the source-local context the proof needs: lines 859--864; lines 874--883.  Nothing was omitted from any retained span, so the package carries no omission record.  No retained line contains a citation, so the wrapper carries no bibliography and no bibliography entry.  No retained line contains a figure environment, an image-inclusion command or drawing code, so no image is delivered and the package declares no asset dependency.  Editorial formatting only, in addition: an AMS \texttt{amsart} preamble that restates the article's own declarations and macros used by the retained lines, namely the environments \texttt{theorem}, \texttt{lemma} on separate counters, together with the standard packages those lines need.  The retained environments print in this document as Theorem 1, Lemma 1 in order of appearance, whereas the article prints the same items with the numbering of its own full document.  The complete native source of the article as distributed in the arXiv e-print for this exact version is bundled unchanged as \texttt{sources/209932/source0.tex}.
\begin{lemma}[Spanning-tree domination]
\label{lem:spanning}
If \(T\) is a spanning tree of a connected graph \(J\), then
\(d_i(J)\leq d_i(T)\) for every \(i\in A\). Consequently,
\(D(J)\leq D(T)\).
\end{lemma}

\begin{lemma}[Rooted-subtree bound]
\label{lem:branch}
Let \(Q\) be a tree of order \(a\), rooted at \(v_0\). For each
\(v\in V(Q)\), let \(t(v)\) be the number of vertices in the subtree rooted at
\(v\), including \(v\). Then
\[
 \sum_{v\in V(Q)}t(v)\leq\frac{a(a+1)}2.
\]
Equality holds if and only if the tree is a path rooted at an endpoint.
\end{lemma}

\begin{theorem}[Sharp aggregate-responsibility bound]
\label{thm:graph-bound}
Let \(A\) be connected with \(m=|A|\geq3\). Then
\begin{equation}
 m\leq R(A)\leq
 F_m:=\left\lfloor\frac{(m+1)^2}{4}\right\rfloor.
 \label{eq:R-bound}
\end{equation}
Equality in the upper bound holds precisely when
\[
 G[A]\cong
 \begin{cases}
 P_3,&m=3,\\
 P_4\text{ or }K_{1,3},&m=4,\\
 P_m,&m\geq5.
 \end{cases}.
\]
Equality in the lower bound holds precisely when deleting each participant
leaves the network induced by the remaining participants connected.
\end{theorem}

\begin{proof}[Proof of Theorem~\ref{thm:graph-bound}]
Because \(H-i\) has \(m-1\) vertices, its largest component has order at most
\(m-1\). Thus \(d_i(H)\geq1\) for every \(i\), and
\[
 R(A)=D(H)\geq m.
\]
Equality holds exactly when \(H-i\) is connected for every \(i\). This proves
the lower bound and its equality statement.

We next prove the upper bound. Lemma~\ref{lem:spanning} reduces the problem to
an arbitrary tree \(T\) of order \(m\). Choose a vertex \(v_0\) such that every
component of \(T-v_0\) has at most \(m/2\) vertices. Every finite tree has such
a vertex. To find one, start at an arbitrary vertex and move into a component of the current
vertex's deletion whenever that component contains more than \(m/2\) vertices.
After a move, the component containing the previous vertex has fewer than
\(m/2\) vertices, and every other component is smaller than the one just
entered. The oversized component therefore shrinks at each step, so the
procedure terminates. This vertex is commonly called a centroid.

Let the components of \(T-v_0\) have orders \(a_1,\ldots,a_p\), and set
\[
 M=\max_j a_j.
\]
Then \(M\leq m/2\) and \(\sum_j a_j=m-1\). Root \(T\) at \(v_0\). For each
nonroot vertex \(v\), let \(t(v)\) be the order of its descendant subtree.
Because \(v\) lies in one branch at \(v_0\), we have
\(t(v)\leq M\leq m/2\). In \(T-v\), the component containing \(v_0\) has
order \(m-t(v)\), while every other component has order at most \(t(v)-1\).
The rootward component is therefore largest, and
\[
 d_v(T)=t(v)\qquad(v\ne v_0).
\]
At the root, the largest component has order \(M\), so
\(d_{v_0}(T)=m-M\). Applying Lemma~\ref{lem:branch} to each branch gives
\begin{equation}
 D(T)\leq m-M+\frac12\sum_{j=1}^p(a_j^2+a_j).
 \label{eq:centroid}
\end{equation}

Suppose first that \(m=2h+1\). Then \(M\leq h\),
\(\sum_j a_j=2h\), and
\[
 \sum_j a_j^2\leq M\sum_j a_j=2hM.
\]
Substitution into \eqref{eq:centroid} yields
\[
 D(T)\leq3h+1+(h-1)M
 \leq3h+1+h(h-1)
 =(h+1)^2=F_m.
\]
For \(h\geq2\), equality requires \(M=h\), and equality in the squared-sum
bound requires every positive \(a_j\) to equal \(h\). Since the branch orders
sum to \(2h\), there are exactly two branches of order \(h\). Equality in
Lemma~\ref{lem:branch} makes both branches paths rooted at the endpoint
adjacent to \(v_0\); hence \(T\cong P_m\). When \(h=1\), the two branches are
singletons and \(T\cong P_3\).

Now suppose \(m=2h\) for an integer \(h\geq2\). If \(M=h\), one branch has order \(h\)
and all remaining branches have total order \(h-1\). Strict superadditivity of
\(g\) bounds their total contribution by \(g(h-1)\), with equality only when
there is one remaining branch. Thus
\[
 D(T)\leq h+\frac{h(h+1)}2+\frac{h(h-1)}2
 =h(h+1)=F_m.
\]
Equality in the rooted-subtree bound makes the two branches endpoint-rooted
paths of orders \(h\) and \(h-1\), so \(T\cong P_m\).

It remains to consider \(M\leq h-1\). Using
\(\sum_j a_j^2\leq M\sum_j a_j=M(2h-1)\) in
\eqref{eq:centroid}, we obtain
\[
 D(T)
 \leq\frac{6h-1+(2h-3)M}{2}
 \leq\frac{2h^2+h+2}{2}.
\]
For \(h\geq3\), this is strictly less than \(h(h+1)=F_m\), since the
difference is \((h-2)/2>0\). If \(h=2\), then \(M\leq1\), and the three
branch orders, whose sum is three, must all equal one. Thus
\(T\cong K_{1,3}\), whose responsibilities are \(3,1,1,1\), and
\(D(T)=6=F_4\). The other equality tree of order four, obtained when \(M=h\),
is \(P_4\).

We have proved the upper bound and classified the equality trees. To verify
attainment, label the vertices of \(P_m\) consecutively. Then
\[
 d_i(P_m)=m-\max\{i-1,m-i\}=\min\{i,m+1-i\}.
\]
Consequently,
\[
 D(P_{2h})=2\sum_{i=1}^h i=h(h+1),
\]
and
\[
 D(P_{2h+1})=2\sum_{i=1}^h i+(h+1)=(h+1)^2.
\]
For \(K_{1,3}\), the center has responsibility three and each leaf has
responsibility one, giving total six.

Finally, we exclude additional edges. Suppose \(H\) attains \(F_m\), and let
\(T\) be an arbitrary spanning tree of \(H\). Spanning-tree domination and the tree
bound imply
\[
 F_m=D(H)\leq D(T)\leq F_m.
\]
Thus \(T\) is one of the equality trees just identified, and
\(d_i(H)=d_i(T)\) for every vertex \(i\).

If \(T\cong P_m\) and \(H\) contains an additional edge, that edge joins two
nonconsecutive path vertices \(v_p\) and \(v_q\). Choose \(v_s\) with
\(p<s<q\). The graph \(T-v_s\) has two nonempty components, whereas the
additional edge connects them in \(H-v_s\). Hence
\(d_{v_s}(H)=1<d_{v_s}(T)\), a contradiction. If
\(T\cong K_{1,3}\), every additional edge joins two leaves. Deleting the
center then leaves a component of at least two vertices in \(H\), whereas all
three components in \(T\) are singletons. The center's withdrawal responsibility
is therefore strictly smaller in \(H\), another contradiction. Hence \(H\)
contains no
edge outside \(T\).

Since \(D(H)=R(A)\), the upper-equality graphs are exactly \(P_3\) for \(m=3\),
\(P_4\) and \(K_{1,3}\) for \(m=4\), and \(P_m\) for \(m\geq5\).
\end{proof}

\end{document}
